\documentclass[a4paper, 12pt]{article}

\usepackage{cmap}
\usepackage[T2A]{fontenc}
\usepackage[english, russian]{babel}
\usepackage[utf8]{inputenc}
\usepackage[left=2cm,right=1.5cm,top=2cm,bottom=2cm]{geometry}
% \usepackage{mathtext}
\usepackage{amsmath}
\usepackage{amssymb}
\usepackage{etoolbox}
\usepackage{amsthm}
\usepackage{nicematrix}
% \usepackage{graphicx}
% \usepackage{tikz}
% \usepackage{parskip}

% \graphicspath{/Images/}

%Реализация aug, overbrace и underbrace без nice matrix
\newcommand\aug{\fboxsep=-\fboxrule\!\!\!\fbox{\strut}\!\!\!}
\newcommand\undermat[2]{\makebox[0pt][l]{$\smash{\underbrace
{\phantom{\begin{matrix}#2\end{matrix}}}_{\text{$#1$}}}$}#2}
\newcommand\overmat[2]{\makebox[0pt][l]{$\smash{\overbrace
{\phantom{\begin{matrix}#2\end{matrix}}}^{\text{$#1$}}}$}#2}
\newcommand\tab[1][.5cm]{\hspace*{#1}}
\newcommand\Underset[2]{\underset{\textstyle #1}{#2}}
\newcommand\Overset[2]{\overset{\textstyle #1}{#2}}


\theoremstyle{definition}
\newtheorem*{definition}{Определение}
\newtheorem*{theorem}{Теорема}
\newtheorem*{consequense}{Следствие}
\newtheorem*{lemma}{Лемма}
\newtheorem*{subtheorem}{Утверждение}
\newtheorem*{remark}{Замечание}


\begin{document}
    \fontsize{14pt}{20pt}\selectfont
    \begin{center}
        \begin{LARGE}
            \textbf{Линейная алгебра и геометрия}
        \end{LARGE}

        \begin{Large}
            \textbf{Глава I. Векторное пространство}
            % \newline
            
            \textbf{\S1. Векторное пространство, размерность, изоморфизм}
        \end{Large}
    \end{center}

    \begin{definition}
        Множество $V$ называется
        \textit{векторным пространством}
        над полем $F$, если заданы операции $+$ и $\cdot$\ :
        $V\times V \to V$, $F \times V \to V$ 
        и выполнены следующие аксиомы: 

        1. $\forall v_1, v_2, v_3\in V$ выполнено
        $(v_1 + v_2) + v_3 = v_1 + (v_2 + v_3)$

        2. $\exists \vec 0 \in V:\ \forall v \in V$
        выполнено $v + \vec 0 = v$

        3. $\forall v \in V \ \ \exists  -v  \in V$:
        $v + (-v) = \vec 0$

        4. $\forall v_1, v_2 \in V$ выполнено 
        $v_1 + v_2 = v_2 + v_1$

        5. $\forall \alpha, \beta \in F, v \in V$ 
        выполнено $(\alpha \beta)v = \alpha (\beta v)$

        6. $\forall v \in V$ выполнено $1 \cdot v = v$

        7. $\forall \alpha, \beta \in F, v \in V$ выполнено
        $(\alpha + \beta)v = \alpha v + \beta v$

        8. $\forall \alpha \in F, v_1, v_2 \in V$ выполнено
        $\alpha (v_1 + v_2) = \alpha v_1 + \alpha v_2$        
    \end{definition}

    \begin{subtheorem}
        Линейным операциям над векторами соотвествует такие же
        операции над их координатами: 
        
        1. $x = eX, y = eY \Rightarrow x + y = e(X + Y)$

        2. $\lambda x = e(\lambda X)$
    \end{subtheorem}
    \begin{proof}
        Пусть $e = (e_1, e_2,...,e_n)$ базис в $V$.
        Тогда $$x = \sum\limits_{i=1}^n x_i e_i =
        (e_1, e_2,..., e_n)
        \begin{pmatrix}
            x_1\\
            x_2\\
            \vdots\\
            x_n
        \end{pmatrix}
        = eX$$
    \end{proof}
    \begin{definition}
        $\varphi: V \to W$ \textit{линейно}, если:

        1. $\forall v_1, v_2 \in V$ выполнено
        $\varphi (v_1 + v_2) = \varphi (v_1) + \varphi (v_2)$

        2. $\forall v \in V, \lambda \in F$ выполнено
        $\varphi (\lambda v) = \lambda \varphi (v)$
    \end{definition}
    \begin{definition}
        $\varphi$ называется \textit{изоморфизмом}, если
        $\varphi$ линейно и биективно.
    \end{definition}
    \begin{consequense}
        Если $\varphi$ линейное отображение, то $\varphi
        (\vec 0_V) = \varphi (\vec 0_W)$
    \end{consequense}
    \begin{proof}
        $$\vec 0_V + \vec 0_V = \vec 0_V \Rightarrow 
        \varphi (\vec 0_V + \vec 0_V) = \varphi (\vec 0_V) +
        (-\varphi (\vec 0_V)) \Rightarrow 
        \varphi(\vec 0_V) = \vec 0_W \Leftrightarrow
        dimV = dimW$$
    \end{proof}
    \begin{consequense}
        $dimV=n \Leftrightarrow V \cong F^n$, где $F^n$
        пространство столбцов высоты n.
        
        ("$\cong$" Чубаров обозначает изоморфность)
    \end{consequense}
    \begin{proof}
        $\Rightarrow$ Пусть $V \cong W$
        изоморфизм $\Rightarrow dimV=dimW$.
        По условию $\exists \varphi: V \to W$
        изоморфизм. Фиксируем базис $(e_1, e_2,..., e_n)$ в $V$ и
        покажем, что \\$\varphi (e_1), \varphi (e_2),...,
        \varphi (e_n)$ базис в $V$.\\
        1. $\varphi (e_1), \varphi (e_2),...,\varphi (e_n)$
        - ЛНЗ и $\sum \limits_{i=1}^{n} \lambda_i \varphi (e_i)
        = \vec 0_W \Rightarrow \sum \limits_{i=1}^{n} \varphi 
        (\lambda_i e_i) = \vec 0_W,$

        $\varphi$ - инъективно, но $\varphi (\vec 0_V) = 0_W
        \Rightarrow \sum \limits_{i=1}^{n} \lambda_i \varphi
        (e_i) = \vec 0_V \Rightarrow$ все $\lambda_i = 0$\\
        2. $\varphi (e_1), \varphi (e_2),..., \varphi (e_n)$
        - полная система: $\forall w \in W\  \exists v \in V:
        \varphi (v) = w$, так как $\varphi$ - сюръективна.
        $\Rightarrow \exists x_1, x_2,..., x_n \in F$. 
        $v = \sum \limits_{i=1}^n x_ie_i \Rightarrow 
        \varphi (v) = w = \sum\limits_{i=1}^n x_i\varphi (e_i)
        \Rightarrow \varphi(e_1), \varphi(e_2),...,
        \varphi(e_n)$ - базис в W.

        $\Leftarrow$ Пусть $dimV = dimW = n.$ Построим
        изоморфизм. Фиксируем базис $e = (e_1, e_2,...,e_n)$
        в $V$ и $f = (f_1, f_2,...,f_n)$ в $W$. Положим:
        $\varphi (e_i) = f_i$. Достроим $\varphi$ до
        линейного отображения, так что $\forall v \in V$
        выполнено $v = \sum\limits_{i=1}^n x_ie_i$.
        Тогда $\varphi (v) = \varphi (\sum\limits_{i=1}^n x_ie_i)
        = \sum\limits_{i=1}^n x_i \varphi (e_i) =
        \sum\limits_{i=1}^n x_if_i$. Таким образом, $\varphi$
        линейно.

    \end{proof}
    \underline{\textbf{Замена базиса}}

    Пусть $dimV = n, e = (e_1,...,e_n)$ - старый базис,
    $e' = (e_1',...,e_n')$ - новый базис.
    
    Пусть известно, что $\forall e_j'$ выполнено $$e_j' = 
    \sum\limits_{i=1}^n c_{ij}e_i = e
    \begin{pmatrix}
        c_{1j}\\
        \vdots\\
        c_{nj}
    \end{pmatrix}
    \Leftrightarrow e'= eC_{e \to e'},$$
    
    где $ C=(c_{ij})$.
    \begin{lemma}
        1. $det\ C\neq 0$

        \ \ \ \ \ \ \ \ \ 2. $C_{e'\to e}=C_{e\to e'}^{-1}$
    \end{lemma}
    \begin{theorem}
        Пусть $e, e'$ два базиса в пространстве $V$,
        $X_e$ и $X_{e'}$ столбцы координат одного и тоже вектора.
        Тогда $X_e = CX_{e'}$
    \end{theorem}
    \begin{proof}
        $\forall x\in V$ выполнено $$x= eX_e= e'X_{e'}= 
        eC_{e\to e'}X_{e'}= e(C_{e\to e'}X_{e'})$$
    \end{proof}
    \newpage
    \begin{center}
        \begin{Large}
            \textbf{\S2. Подпространство}
        \end{Large}
    \end{center}
    \begin{definition}
        Подмножество $U$ в пространстве $V$ называется\\
        \textit{подпространством} в $V$, если:

        1. $\forall U\neq 0$

        2. $\forall u_1, u_2 \in U$ выполнено $u_1 + u_2 \in U$

        3. $\forall u \in U, \lambda \in F$ выполнено
        $\lambda u\in U$
    \end{definition}
    \textbf{Способы задания подпространства:}

    1. $U =\langle a_1, a_2,..., a_n\rangle$

    2. $dimV, W = \{v = eX|AX = 0\}$

    \begin{definition}
        \textit{Линейная оболочка векторов} $a_1, a_2,...,a_n
        \in V$ -- это\\ $\{\lambda_1a_1 + \lambda_2a_ + ... + 
        \lambda_na_1n|\lambda_i \in F\}$
    \end{definition}
    \begin{lemma}
        1. $\langle a_1, a_2,..., a_m \rangle := U$ -- подпространство в $V$

        2. $dim\ U = rk\{a_1, a_2,..., a_m\}$
    \end{lemma}
    \begin{proof}
        Если $dim\ U = n$, то в фиксируем базис и составляем
        матрицу из столбоцв координат $a_1, a_2,..., a_m:$
        $$A = (a_1^\uparrow, a_2^\uparrow,...,a_m^\uparrow)\sim
        \begin{pmatrix}
            a_{1j_1} & \hdotsfor{7}\\
            0 & a_{1j_2} & \hdotsfor{6}\\
            \hdotsfor{8}\\
            0 & 0 & \dots & \hdotsfor{3} & a_{rj_r} & \dots\\
            \hdotsfor{8}\\
            0 & 0 & \hdotsfor{6}
        \end{pmatrix}$$
        Элементарными преобразованиями строк приводим матрицу A
        к ступенчатому виду и получаем, что столбцы $a_{j_1},
        a_{j_2},..., a_{j_r}$ составляют бызис в $U$.\\
    \end{proof}
    \begin{lemma}
        Любую ЛНЗ систему векторов в $V (dimV < \infty)$ можно
        можно дополнить до басима пространства $V$.
    \end{lemma}
    \underline{\textbf{Алгоритм}}: Пусть $a_1, a_2,..., a_m$ - ЛНЗ,
    $m < n = dimV$, известны координаты этих векторов в
    некотором базисе. Тогда $rk(a_1^\uparrow, a_2^\uparrow,...,
    a_m^\uparrow) = m.$ Составим матрицу столбцов из них и
    припишем столбцы $E$ порядка n:
    $$(A|E_n) \sim 
    \begin{pmatrix}
        \text{Выделим базисные столбцы}\\
        \text{Э.П. строк в \(A|E_n\),}\\
        \text{включая столбцы A}
    \end{pmatrix}
    \sim rk(A|E_n) = n$$\\
    Вывод: к $a_1, a_2,..., a_m$ надо добавить единичные столбцы,
    вошедшие в базис матрицы $A|E_n$.
    \newpage
    \begin{center}
        \underline{\textbf{Операции с блочными матрицами}}
    \end{center}
    \begin{definition}
        \textit{Блочная матрица} -- матрица, разбитая на
        подматрицы, которые обозначаются отдельными буквами.
    \end{definition}
    \textbf{Пример}: $$1.\ A = 
    \begin{pmatrix}
        3 & -1 & \aug & 2 & 0\\
        2 &\ 1 & \aug & -1 & 4
    \end{pmatrix}
    = (A_1 | A_2) \text{ - блочная строка.}$$

    $$2.\ B = 
    \begin{pmatrix}
        1 & 0\\
        0 & -1\\
        \hline
        1 & 3\\
        4 & -2
    \end{pmatrix} = 
    \begin{pmatrix}
        B_1\\
        -\\
        B_2
    \end{pmatrix} \text{ - блочный столбец.}$$
    1. Линейная оболочка.\\ 2. ОСЛУ.\\
    $2 \Rightarrow 1:$ строим ФСР. $AX = 0$. Проводим элементарные
    преобразования строк, перенумеровывая неизвестные, таким
    образом чтобы привести матрицу к следующему виду:
    
    $$A \underset{\text{строк}}{\overset{\text{Э.П.}}{\sim}}
    \begin{pmatrix}
        \overmat{r}{1 & \dots & 0} & a_{1j} & \dots\\
        \vdots & \ddots & \vdots & \vdots & \vdots\\
        0 & \dots & 1 & a_{rj} & \dots\\
        % 0 & \dots & 0 & 0 & \dots\\
        \hdotsfor{5}\\
        0 & \dots & 0 & 0 & \dots
    \end{pmatrix}
    , r + 1 \leq j \leq n$$
    Возвращаемся к уравнениям: $x_i = -\sum\limits_{j = r + 1}^
    n a_{ij}x_j$, где $1\leqslant i \leqslant r$, а
    ($x_{r+1},...,x_n$) -- свободные неизвестные.
    В векторном виде: $$X =
    \begin{pmatrix}
        x_1\\\vdots\\x_r\\x_{r+1}\\\vdots\\x_n
    \end{pmatrix}
    =
    \begin{pmatrix}
        -\sum\limits_{j=r+1}^na_{1j}x_j\\\vdots\\
        -\sum a_{rj}x_j\\x_{r+1}\\
        \vdots\\x_n
    \end{pmatrix}
    = x_{r+1}
    \undermat{Y_1}{\begin{pmatrix}
        -a_{1,r+1}\\\vdots\\-a_{r,r+1}\\1\\\vdots\\
    \end{pmatrix}}
    +...+ x_n
    \undermat{Y_n}{\begin{pmatrix}
        -a_{1,r+1}\\\vdots\\-a_{r,r+1}\\0\\\vdots\\
    \end{pmatrix}}$$ где $Y_1,..., Y_n$ -- фундаметальная
    система решений, т.е. базис пространства решений $AX = 0$.
    \newpage
    Составим матрицу из столбцов ФСР.
    $$\varPhi = 
    \begin{pmatrix}
        -a_{1,r_1} & \dots & -a_{a,n}\\
        \vdots & \null & \vdots\\
        1 & \null & 0\\
        % 0 & \ddots & \vdots\\
        \vdots & \ddots & \vdots\\
        0 & \null & 1
    \end{pmatrix}
    =
    \begin{pmatrix}
        -B\\\hline E_{n-r}
    \end{pmatrix}
    - \text{\ фундаме нтальная матрица.}$$
    В общем случае фундаметальная матрица -- это матрица,
    имеющая $n = rkA$ ЛНЗ столбцов, которые являются
    решениями системы $AX = 0$, где $A \sim (E_r|B)$
    \begin{theorem}
        Способы 1 и 2 равносильны, если $dinV<\infty$
    \end{theorem}
    \begin{proof}
        $1 \Rightarrow 2$. Даны ЛНЗ векторы $c_1, c_2,...,c_m$
        в $V$.\\ Нужно найти такую матрицу $A_{p\times n}$,
        чтобы $W = \{x \in F^n|AX = 0\}$, где $X$ - столбец
        координат произвольного вектора из $U = \langle c_1,c_2,..., c_m \rangle$
        Если составить матрицу из столбцов $c_1, c_2,..., c_m$,
        то оно должна стать фундаметальной матрицей для ОСЛУ: $AX = 0$.
        \\\underline{1 этап:} Векторы $c_1, c_2,...,c_m$
        расписать по строкам:
        $$\begin{pmatrix}
            \vec c_1\\\vdots\\\vec c_m
        \end{pmatrix}\underset{\text{строк}}{\overset{\text{Э.П.}}{\sim}}
        (E_n|\mbox{\scriptsize{\undermat{{\mbox{\normalsize m}}}
        {\null}}}B)$$
        \underline{2 этап:} $$\varPhi = 
        \begin{pmatrix}
            B^T\\
            \hline E_m
        \end{pmatrix} \rightsquigarrow A = (E_{n-m}|-B^T)
        \text{ -- искомая матрица системы.}$$
        Значит, $p = n - m.$
    \end{proof}
    \begin{remark}
        Столбец $\begin{pmatrix}x_1\\\vdots\\x_n\end{pmatrix}$
        принадлежит $c_1,...,c_m \Leftrightarrow
        rk\Biggl(c_1^\uparrow,...,c_m^\uparrow |
        \begin{pmatrix}x_1\\\vdots\\x_n\end{pmatrix}\Biggr) =
        rk(c_1^\uparrow,...,c_m^\uparrow)$
    \end{remark}
    \begin{proof}
        Составить матрицу:$$
        rk\Biggl(c_1^\uparrow,...,c_m^\uparrow |
        \begin{pmatrix}x_1\\\vdots\\x_n\end{pmatrix}\Biggr)
        \underset{\text{строк} }{\overset{\text{Э.П.}}{\sim}}
        \begin{pmatrix}
            c_{11} & \null & \null & \vline & \hdotsfor{1}\\
            \null & \ddots & \null & \vline & \vdots\\
            \null & \null & c_{mm} & \vline & \hdotsfor{1}\\
            \hline
            0 & \dots & 0 & \vline & \sum a_{m+1,j} x_j\\
            \vdots & \ddots & \vdots & \vline & \vdots\\
            0 & \dots & 0 & \vline & \sum a_{n,j} x_j\\
        \end{pmatrix},$$ где суммы $\sum a_{i,j}x_j := 0,
        \text{где} i = m + 1,...,n.$
    \end{proof}
    \newpage
    \begin{center}\begin{Large}
            \textbf{\S3. Пересечение и сумма подпространств.}
    \end{Large}\end{center}
    \begin{subtheorem}
        Пусть $U_i$ -- семейство подпространств векторного
        пространства $V$. Тогда $\bigcap\limits_{i \in I} U_i$
        -- подпространство $V$.
    \end{subtheorem}
    \begin{proof}
        $\\1.\ \vec 0_V \in U_i,\ \forall i \in I \Rightarrow
        \vec 0_V \in \bigcap\limits_{i \in I} U_i$\\
        $2. \forall x, y \in U_i,\ \forall i \in I \Rightarrow
        x + y \in U_i, \forall i \in I \Rightarrow
        (x + y) \in \bigcap\limits_{i \in I} U_i$\\
        $3. \forall x \in U_i,\ \forall i \in I,\ \forall \lambda
        \in F \Rightarrow \lambda x \in U_i, \forall i \in I
        \Rightarrow \lambda x \in \bigcap\limits_{i \in I}U_i$\\
        $U_1 \bigcup U_2$ может быть подпространством:
        $u_1 + u_2 \notin U_1 \bigcup U_2$\\
    \end{proof}
    \begin{definition}
        \textit{Суммой подпространств} $U_1,...,U_m \subset V$ называется
        множество $U_1 +...+ U_m = \{u_1 +...+ u_m|u_i \in U_i\}.$
    \end{definition}
    \begin{subtheorem}
        $U_1...+ U_m$ -- подпространствов $V, dim(U_1 +...+ U_m)
        \leqslant \sum\limits_{i=n}^n dim\ U_i$
    \end{subtheorem}
    \begin{proof}
        Если $e_1,..., e_m$ -- базисы в этих подпространствах,
        то\\ $U_1 +...+ U_m = \langle e_1,..., e_m \rangle:\ U_1 = \sum
        \limits_{i=1}^{n_1} \lambda_{1i}e_{1i},...,
        U_m = \sum\limits_{j=1}^{n_m} \lambda_{mj}e_{mj}
        \Rightarrow\\ U_1 +...+ U_m$ -- линейная комбинация этих
        векторов.\\
    \end{proof}
    \underline{\textbf{Формула Грассмана.}} Если $U_1, U_2$ --
    конечные подпространства $V$, то $dim(U_1 + U_2) = dim\ U_1
    + dim\ U_2 - dim(U_1 \cap U_2)$
    \begin{proof}
        $dim\ U_1 = n_1,\ dim\ U_2 = n_2,\ dim(U_1 \cap U_2) = r.$
        \\Выберем $e_1,..., e_r$ -- базис в $U_1 \cap U_2$.
        Дополним их до базиса в $U_1$ векторами $a_{r+1},...,a_{n_1},$
        в $U_2$ векторами $b_{r+1},...,b_{n_2}$    
    \end{proof}
    \begin{subtheorem}
        $\{c_i, a_j, b_k\}$ -- базис $U_1 + U_2$ (их количество
        $n_1 + n_2 - r$).
    \end{subtheorem}
    \begin{proof}
        Ясно, что это полная система.
        Докажем ЛНЗ:\\ Допустим $\sum\limits_{i=r+1}^{n_1}
        \alpha_i a_i + \sum\limits_{j=r+1}^{n_2}\beta_j b_j +
        \sum\limits_{k=1}^r\gamma_k c_k = 0\ \Longrightarrow$
        \begin{flushright}
            $\undermat{\in U_1}{\sum\alpha_i a_i + \sum\gamma_k
            c_k} = \undermat{ \in U_1 \cap U_2}{-\sum\beta_j b_j}
            \Longrightarrow$            
        \end{flushright}
        $\\\\\gamma_k':= -\sum\beta_jb_j = \sum\gamma_r'c_r\
        \Longrightarrow \sum\beta_jb_j + \sum\gamma_k'c_k = 0 
        \stackrel{\text{ЛНЗ}}{\ \Longrightarrow\ } \beta_j = 0,
        \forall j = r + 1,..., n_2;\\ \gamma_k' = 0 \Longrightarrow
        \sum\alpha_ia_i + \sum\gamma_kc_k = 0 
        \stackrel{\text{ЛНЗ}}{\ \Longrightarrow\ }
        \alpha_i = 0,\ \forall i = r+1,..., n;
        \gamma_k = 0,\ \forall k = 1,..., r.$    
    \end{proof}
    \newpage
    \underline{\textbf{Алгоритм вычисления базисов}} в
    $U_1 + U_2,\ U_1 \cap U_2\ (dimV \leqslant \infty).$
    \begin{remark}
        $U_1 + U_2 = \langle a_1,..., a_{n_1}, b_1,..., b_{n_2} \rangle$
    \end{remark}
    \begin{proof}
        Можно составить матрицу из столбцов координат и\\ Э.П.
        строк выявить базисные столбцы в этой расширенной матрице.
        \\Вектор $v \in U_1 \cap U_2 \Longleftrightarrow
        \exists\ x_i, y_i \in F:\ v = \sum\limits_{i=1}^{n_1}
        x_ia_i =\ \sum\limits_{j = 1}^{n_2}y_jb_j$, т.е.\\
        $(x_1,..., x_{n_1}, -y_1,..., -y_{n_2})$ -- решение
        ОСЛУ с той же самой матрицей.\\ Базис $U_1 \cap U_2$
        будет давать ФСР (ее часть соответственно $\{a_i\}$
        или $\{b_i\}$).\\
    \end{proof}
    1. Составить матрицу:
    $$(\undermat{\text{\scriptsize ЛНЗ}}{a_1^\uparrow... a_{n_1}^\uparrow}|
    \ \undermat{\text{\scriptsize ЛНЗ}}{b_1^\uparrow... b_{n_2}^\uparrow})
    \underset{\text{строк}}{\overset{\text{Э.П.}}{\sim}}
    \begin{pmatrix}
        \overmat{E_{n_1}}{
        1 & \dots & 0} & \vline & 0 & \dots & 0 & \vdots & \vdots\\
        \vdots & \ddots & \vdots & \vline & \vdots & \null & \vdots & \vdots & \vdots\\
        0 & \dots & 1 & \vline & 0 & \dots & 0 & \vdots & \vdots\\
        \hline
        0 & \dots & 0 & \vline & 1 & \dots & 0 & \vdots & \vdots\\
        \vdots & \null & \vdots & \vline & \vdots & \ddots & \vdots & \vdots & \vdots\\
        0 & \dots & 0 & \vline & 0 & \dots & 1 & \vdots & \vdots\\
    \end{pmatrix}$$
    Элементарными преобразованиями строк приводим к улучшенному
    ступенчатому виду. Ветокторы-столбцы $a_1,..., a_{n_1},
    b_1,...,b_{n_2}$ -- базис $U_1 + U_2$.

    2. Вычислить $dim(U_1 \cap U_2) = n_1 + n_2 - m$ (это
    количество столбцов $b_j$ не вошедших в базис суммы).
    $m = n_1 + k = dim(U_1 + U_2).$

    3. Выразить векторы $b_{k+1},...,b_{n_2}$ через базисы:
    $b_l = \sum\limits^m\alpha_{il}a_i + \sum\limits^k\beta_{sl}b_s$
    -- ЛНЗ $\Longleftrightarrow b_l - \sum\beta_{sl}b_s =
    \sum\alpha_{il}a_i \in U_1 \cap U_2$ (Их количество $dim(U_1 \cap U_2)$).
    \begin{center}
        \begin{Large}
            \textbf{\\\S4. Прямая сумма}
        \end{Large}
    \end{center}
    Пусть $V$ -- векторное пространство над полем $F$,\
    $U_1,..., U_k$ -- подпр-ва в $V$.
    \begin{definition}
        Сумма $U_1 +...+ U_k \ (k \geqslant 2)$ называется
        \textit{прямой суммой пространств} $U_1,..., U_k$, если
        $\forall u$ из суммы представляется в виде $u = u_1 +...+u_k$
        единственным образом. 
    \end{definition}
    \underline{Обозначение:} $U_1\oplus...\oplus U_k$\\
    \textbf{Примеры:} \\1. $U_1 = \{A \in M_n(\mathbb{R})|A^T = A\}$
    -- симметричные матрицы\\
    2. $U_1 = \{B \in M_n(\mathbb{R})|B^T = B\}$
    -- кососимметричные матрицы
    \newpage
    \begin{subtheorem}
        $M_n(\mathbb{R}) = U_1 \oplus U_2$, где\ \ $U_1$ --
        пространство симметричных\\ матриц, а $U_2$ -- 
        пространство кососимметричная матриц.
    \end{subtheorem}
    \begin{theorem}
        Следующие условия равносильны:

        $1.\ U_1 + U_2 = U_1 \oplus U_2$

        $2.\ U_1 \cap U_2 = \{0\}$

        $3.\ dim(U_1 + U_2) = dim\ U_1 + dim\ U_2$

        4. Базис в $U_1 + U_2$ -- объединение базисов слагаемых.
    \end{theorem}
    \begin{proof}
        $\\1 \Rightarrow 2:\ \forall u \in U_1 + U_2$ выполнено
        $u = u_1 + u_2$ -- единственным образом.
        Допустим противное: пусть $\exists\ u_0 \in U_1 + U_2,\ 
        u_0 \neq 0 \Longrightarrow u_0 = \mbox{\scriptsize
        \undermat{\in U_1}{\null}}u_0 + \mbox{\scriptsize
        \undermat{\in U_2}{\null}} 0 = \mbox{\scriptsize
        \undermat{\in U_1}{\null}} 0 + \mbox{\scriptsize
        \undermat{\in U_2}{\null}} u_0$.
        Противоречие единственности $\Longrightarrow U_1 \cap U_2
         = \{0\}.$\\
        $2 \Rightarrow 3:$ По формуле Грассмана.\\
        $3 \Rightarrow 4:$ Ясно, что $U_1 + U_2 =
        \langle a_1,...,a_{n_1}, b_1,..., b_{n_2} \rangle$ -- эти векторы ЛНЗ,
        если $$\exists \sum_{i=1}^{n_1}\alpha_ia_i +
        \sum_{j=1}^{n_2}\beta_jb_j = 0 \Longrightarrow
        \undermat{= 0}{\sum\alpha_ia_i} = -\undermat{=0}
        {\sum\beta_jb_j} = \{0\}\in U_1 \cap U_2$$\\
        $4 \Rightarrow 1:\forall u = u_1 + u_2 =
        \sum\limits_{i=1}^{n_1} x_ia_i + \sum\limits_{j=1}^{n_2} y_jb_j$
        раскладывается по базису единственным образом $\Rightarrow
        u_1, u_2$ единственны.
    \end{proof}
    \begin{theorem}
        Для $U_1,..., U_k\ (k \geqslant 2)$ следующие условия
        равносильны:

        $1.\ U_1+...+U_k = U_1 \oplus...\oplus U_k$

        $2.\ \forall i = 1,...,k \text{ выполнено } U_i\cap(
        \sum\limits_{j \neq i}U_j) = \{0\}$

        $3.\ dim(U_1+...+U_k) = dim\ U_1+...+dim\ U_k$

        4. Базис в $U_1+...+U_k$ -- объединение базисов слагаемых.
    \end{theorem}
    \begin{subtheorem}
        Для любого пространства $U \subset V\ \exists \text{
        подпространство } W \subset V: V = U \oplus W$,
        где $W$ -- прямое дополнение к $U$.
    \end{subtheorem}
    \begin{proof}
        Пусть $e_1,..., e_k$ -- базис в $U$, тогда $\exists$
        векторы $e_{k+1},...,e_n\\ (n = dimV): \{e_1,...,e_n\}$ --
        базис в $V$, тогда $W := \langle e_{k+1},...,e_n \rangle$
    \end{proof}
    \newpage
    \begin{center}\underline{\textbf{Факторпространство}}\end{center}
    Пусть $V$ -- векторное пространство, $U \subseteq V$ --
    подпространство.\\ Скажем, что векторы $v_1$ и $v_2 \in V$
    сравнимы по модулю, если $v_1 - v_2 \in U$, то есть
    $v_1 \equiv v_2\ (mod\ U)$.
    \begin{subtheorem}
        Отношение: $v_1 \sim v_2 \Longleftrightarrow u_1 - u_2
        \in U$ является отношением эквивалентности на $V$.
    \end{subtheorem}
    \underline{Класс эквивалентности вектора $v$:}\\
    $\vec v = \{v + u\ |\ \forall u \in U\} = v + U\\
    V = \bigsqcup\limits_{v \in V}(v + V)$\\
    Сравним с видом решения системы $AX = b: X = X_r + Y$ -- 
    частное решение + общее решение однородной ассоциированной
    $AY = 0$.  Классы эквивалентности множества решений $AX = b$.
    \\\underline{Обозначим:} $V/U$ -- множество классов
    эквивалентности (смежных классов по $U$).\\
    \underline{Термин:} $V/U$ -- факторпространство $V$ по $U$\\
    \underline{Определим сложение классов:} $(v_1 + U) + (v_2 + U
    := v_1 + v_2 + U)$\\
    \underline{Определим умножение классов на скаляр $\lambda
    \in F:$} $\lambda(v+U) := \lambda v + U$
    \begin{subtheorem}
        Множество с введенными операциями является векторным
        \\пространством.
    \end{subtheorem}
    \begin{subtheorem}
        1. Если $dimV < \infty,$ то $dim(V/U) = dimV - dim\ U$

        \ \ \ \ \ \ \ \ \ \ \ \ \ \ \ \ \ \ \ \ 
        2. $(U\oplus V) / U \cong W$
    \end{subtheorem}
    \begin{proof}
        $\\2.$ I способ: Построим отображение: $\forall v \in V\
        \exists!\ u, w, \text{ такие что }\\ f: v = u + w \mapsto v + U,\ f$ -- линейная.
        $f$ -- сюръективна: $w + u = f(w).\\ f$ -- инъективна,
        так как если $w_1 + u = w_2 + u \Rightarrow w_1 - w_2 \in U.$
        \\II способ: $W$ имеет единственный общий вектор с любым
        смежным классом по $U$.
        $\forall x \in V = U + W\ \exists!\ u \in U,\ v \in V
        \text{:}\ 
        v = u + w \Longrightarrow$ Рассмотрим $v + u = w + (u + U) =
        w + U$. Построим отображение $\varphi: V/U \mapsto W$
        по правилу $\varphi(w + U) =  w$, причем $\varphi$ --
        линейное.\\ Тогда $\varphi((w_1 + U) + (w_2 + U)) =
        \varphi(w_1 + w_2 + U) = w_1 + w_2 = \varphi(w_1 + U) + 
        \varphi(w_2 + U)$.\ $\forall \lambda \in F, \varphi(
        \lambda w + U) = \lambda w = \lambda \varphi(w + U).\ $ 
        $\varphi$ -- биективное, $\forall w \in W$,\\ то $\varphi(
        w + U) = w$,
        если $\varphi(w + U) = 0 \in W \Longrightarrow
        w + U = U \Longrightarrow v = 0 \Longrightarrow\\ 0$ -- единственный.
        \newpage
        1. Рассмотрим базис $V$, составленный из базисов слагаемых:\\
        $\undermat{\text{Базис в U}}{e_1,...,e_{n_1}},
        \undermat{\text{Базис в W}}{e_{n_1 + 1},...,e_n}\ (n = dimV)$
        $$\forall v \in V: v = \sum\limits_{i = 1}^{n_1}x_ie_i + 
        \sum\limits_{j = n_1 + 1}^nx_je_j \Rightarrow \overline{v} =
        \sum\limits_{j = n_1 + 1}^nx_j\overline{e_j}$$
        Т.е. смежные классы векторов $e_{n_1 + 1},..., e_n$ --
        полная система в $V/U.$\\ Они ЛНЗ если $\sum\limits_{
        j = n_1 + 1}^nx_j\overline{e_j} = \overline{0} = U
        \Longrightarrow \sum\limits_{j = n_1 + 1}^nx_je_j
         \in U \cap W = \{0\} \Longrightarrow$
        Все $\lambda_j = 0 \Longrightarrow \{\overline{e}_{n_1 + 1},
        ...,\overline{e_n}\}$
        -- базис $V/U.$
    \end{proof}
    \begin{remark}
        Если $v = 
        \begin{pmatrix}
            x_1\\\vdots\\x_{n_1}\\x_{n_1 + 1}\\\vdots\\x_n
        \end{pmatrix}$
        , то $\overline{v} = 
        \begin{pmatrix}
            x_{n_1 + 1}\\\vdots\\x_n
        \end{pmatrix}$
    \end{remark}
    \underline{\textbf{Внешняя прямая сумма пространств:}}
    Пусть $V_1,..., V_m$ -- векторные пространства над полями $F$.
    \\\underline{Обозначение:} $V = V_1 +...+ V_m =
    \{(v_1,..., v_m)|v_i \in V_i, 1 \leq i \leq m\}$, $V$-- 
    внешняя прямая сумма пространств.
    \begin{remark}
        Пространство $V = V_1 +...+ V_m$ можно превратить в
        прямую сумму подпространств: рассмотрим $U_i = \{
        (0_{V_1},..., v_i,..., 0_{V_m})|v_i \in V_i\} \subset V.
        \\(v_1,...,v_m) = (v_1, 0,...,0) + (0, v_2, 0,...,0) +...
        \Longrightarrow V = U_1 \oplus...\oplus U_m.$
    \end{remark}
    В факторпространстве $V/U: \underline{0} = U$ и $
    -\overline{v} = \overline{(-v)}$

    \begin{definition}
        $dim(V/U)$ -- коразмерность пространства $U$ в
        пространстве $V$. Если $dim V < \infty, \text{ то }
        codim_{V}(U) = dimV - dim\ U.$
    \end{definition}
    \newpage
    \begin{center}\begin{Large}
        \textbf{\S5. Линейные функции и сопряженное (двойственное)
        пространство.}
    \end{Large}\end{center}
    Пусть $V$ -- векторное пространство над полем $F$.
    \begin{definition}
        Функция $f: V \mapsto F$ называется \textit{линейной}
        если:
        
        $1.\ \forall v_1,v_2 \in V \Longrightarrow f(v_1 + v_2)
        = f(v_1) + f(v_2)$

        $2.\ \forall v \in V, \lambda \in F \longrightarrow
        f(\lambda v) = \lambda f(v)$
    \end{definition}
    \begin{definition}
        \textit{Ядро функции:} $Ker f = \{v \in V|f(v) = 0\}$
    \end{definition}
    $\\$Образ функции $Imf = f(V) = \{\alpha \in F|\ \exists v \in V:
    f(v) = \alpha\}$
    \begin{subtheorem}
        $Kerf$ -- подпространство в $V^*$. Если $f \not\equiv 0$,
        то $Kerf$ имеет коразмерность $1$ в $V$.
    \end{subtheorem}
    \begin{proof}
        $$\forall v_1, v_2 \in Ker f, \lambda, \mu \in F
        \Longrightarrow f(\lambda v_1 + \mu v_2) =
        \lambda f(v_1) + \mu f(v_2)  = 0$$
    \end{proof}
    \underline{Обозначение:} $V^* = \{f: V \mapsto F\}$ --
    линейная функция.
    
    Определим на $V^*$ операции сложения: $\forall f_1, f_2 \in V^*$
    выполнено $(f_1 + f_2)(v) := f_1(v) + f_2(v)$. Умножение на
    элемент поля: $\forall f \in V^*, \lambda \in F$ выполнено
    $(\lambda f)(v) := \lambda f(v)$

    \underline{Термин:} $V^*$ -- пространство сопряженное к $V$.
    \begin{theorem}
        Если $dimV = n$, то $dimV^* = n$ (следовательно $V^* = V$).
    \end{theorem}
    \begin{proof}
        Введем в $V$ базис $e_1,..., e_m, \forall v \in V:
        v = \sum\limits_{i=1}^nx_ie_i.$\\ Тогда $\forall f
        \in V^*,\ f(v) = \sum\limits_{i=1}^nx_if(e_i) =
        \sum\limits_{i=1}^nf(e_i)f_i(v) = (\sum\limits_{i=1}^n
        f(e_i)f_i)(v),\forall v\in V.$
        $(f(e_1),..., f(e_n))$ -- строка коэффициентов функций $f$.
        $f = \sum\limits_{i=1}^nf(e_i)f_i(v) = \sum\limits_{i=1}^n
        a_if_i(v).$ Таким образом $V^* = \langle f_1,..., f_n \rangle$.
        Они ЛНЗ, если $\exists \lambda_1,..., \lambda \in F:
        \undermat{\widetilde{f}}{\sum\limits_{i=1}^n\lambda_if_i = 0}.$
        $\widetilde{f}(e_j) = \sum\limits_{i=1}^n\lambda_if_i(e_j)
         = \lambda_i = 0,\ 1 \leq j\leq n$.\ 
        $f_i(e_j) = \begin{cases}
            1, j = i\\0, j \neq i
        \end{cases} \Longrightarrow\\\\\\ \lambda_1=...=\lambda_n=0.$\\
    \end{proof}
    \newpage
    \textbf{Примеры:}

    $1. V = \mathbb{R}[x], x_0 \in \mathbb{R}$ -- фиксированная
    точка. $f: p(x) \mapsto p(x_0).$\\
    Рассмотрим $V_n = \{p(x)|deg\ p \neq n\}, \forall n \in \mathbb{N}.$
    \\Определим $Ker f = \{p(x)\ |\ p(x_0) = 0\}$. Базис:
    $(x - x_0), (x-x_0)^2,...,(x-x_0)^n$.

    $2. dimV = n, \text{ пусть } e_1,..., e_n$ -- базис в $V.\
    \forall v \in V, v = \sum\limits_{i=1}^nx_ie_i$.
    Рассмотрим $f_i(v) = x_i$ -- $i$-ая координатная функция.
    \begin{definition}
        Обозначим $f_i := e^i,\ 1 \leq i \leq n.\ \{e^1,...,e^n\}$
        -- базис в $V^*$, двойственный (биортогональный) к базису
        $e_1,..., e_n$.
    \end{definition}
    По построению $e^i(e_j) =
    \begin{cases}1, j = i\\0, j \neq i\end{cases}$
    (символ Кронекера)

    Разложение $\forall v \in V$ по базису $e_1,..., e_n$ имеет
    вид $v = \sum\limits_{i=1}^ne^i(v)\ e_i$
    \begin{subtheorem}
        Строки коэффициентов линейных функций: $f(v) =
        \sum\limits_{i=1}^na_ix_i$ изменяется при переходе к
        новому базису по формулам: $\vec a = (a_1,..., a_n),\\
        \vec{a'} = \vec a C_{e\to e'}$
    \end{subtheorem}
    \begin{proof}
        $f(v) = \vec{a'}x'^{\uparrow}$. Знаем, что $X = CX'
        \Longleftrightarrow \vec a(CX') = \vec{a'}X'
        \Longleftrightarrow (\vec a\ C)X' = \vec{a'}X'$, верно
        $\forall x' \in F^n.$
    \end{proof}
    \textbf{Примеры взаимных базисов:}\\Пусть $V = \mathbb{R}_n
    [x] = \{p(x) = a_0 + a_1x +...+ a_nx^n|\ n \in \mathbb{N},\ 
    a_i \in \mathbb{R}\}$

    1. С базисом: $1, \frac{x - x_0}{1!},..., \frac{(x-x_0)^n}{k!},
    \ x_0 \in \mathbb{R}$ s-- фиксированное число.
    $$p(x) = \sum\limits_{k=0}^n \frac{p^{(k)}(x_0)}{k!}(x-x_0)^k$$
    Можно рассмотреть линейные функции $\delta^{(k)}: \delta^{(k)}
    (p) := p^{(k)}(x_0),\ k = 0,...,n.\\
    \{\delta^{(k)}, k = 0,...,n\}$ -- базис в $V^*$ двойственный
    к тейлоровскому базису.

    2. $V = \mathbb{R}_n[x]$. Пусть $x_0,...,x_n$ -- попарно
    различные числа. Рассмотрим\\ линейные функции
    $\varphi_k(p) = p(x_k),\ k = 0,.., n.$
    $$\forall p(x) = \sum\limits_{k=0}^np(x_k)l_k(x),\ \ l_k(x_i)
    = \delta_{ik},\ \ l_k(x) = \frac{\prod_{i\neq k}(x - y_i)}
    {\prod_{i\neq k}(x_k - x_i)}$$

    \underline{Обозначение:} $V^{**} = (V^*)^*$ -- второе сопряженное
    \newpage
    \begin{theorem}
        Если $dimV < \infty$, то $V^{**} \cong V$, причем
        изоморфизм не зависит от базиса.
    \end{theorem}
    \begin{proof}
        Построим отображение $\varphi:\ V \mapsto V^{**},\
        \forall v \in V:\\ \varphi(v) := \varphi_v,\ \varphi_v
        \in V^{**}.\ \ \varphi_v(f) := f(v),\ \forall f \in V^{**}.$
        Ясно что $\varphi$ -- линейное отображение: $\forall f
        \in V^*\ \ \varphi_{v_1 + v_2}(f) = f(v_1 + v_2) =
        \varphi_{v_1}(f) + \varphi_{v_2}(f)$, и $\forall \lambda
        \in F\\ \varphi_{\lambda v}(v) = f(\lambda v) =
        \lambda f(v) = \lambda \varphi_v(f)$\\
        $\varphi$ -- инъективное отображение $\Longleftrightarrow
        Ker\varphi = 0,\\ Ker\varphi = \{v \in V|\ f(v) = 0,\
        \forall f \in V^*\} \{0\}$. Таким образом, $\varphi:\ 
        V \longmapsto W$ --\\ инъективное линейное отображение. Но
        $dimV = dim(V^*)^* \Longrightarrow \varphi$ --
        сюръективное отображение.\\\textit{Подробнее:} если
        $\varphi:\ V\longmapsto W$ -- инъективное линейное отображение
        и $dimV = dimW$, то оно сюръективно: $\exists e_1,...,e_n$
        -- базис в $V$, тогда $\varphi(e_1),..., \varphi(e_n)$--
        ЛНЗ $\Longrightarrow \varphi(e_1),..., \varphi(e_n)$ --
        базис в $W$. Рассмотрим $\lambda_1\varphi(e_1)+...+
        \lambda_n\varphi(e_n) = 0_W \Longleftrightarrow
        \varphi(\lambda_1 e_1+...+ \lambda_n e_n) = 0_W$.
        Но $\varphi(0_V) = 0_W.$\\ Инъективность $\Longrightarrow
        \undermat{\text{ЛНЗ}}{\lambda_1 e_1+...+ \lambda_n e_n}
        = 0_V \Longrightarrow \lambda_1=...=\lambda_n = 0$.\\
    \end{proof}
    \begin{theorem}
        Векторы $a_1,...,a_k \in V\text{-- ЛНЗ}\ (dimV < \infty)
        \exists f_1,..,f_k \in V^*:\\ det(f_i(a_j)) \neq 0$.
    \end{theorem}
    \begin{proof}
        $\\\Longrightarrow$ Дополним векторы до базиса $V$:
        $a_1,..., a_k,..., a_n$\\ По нему строим $f_1,...,f_n$
        -- двойственный базис. Тогда $f_j(a_i) = \delta_{ij}.$
        по определению $f_j(a_i) = E \Longrightarrow det(f_j(a_i))
        \neq 0$\\$\Longleftarrow$ Пусть $P := (f_j(a_i)),\ 
        P \in Mat_{k\times k}$. Предположим, что $P$ --
        невырожденная, тогда $(a_1',..., a_k') = (a_1,..., a_k)
        P^{-1}$\\$(f_j(a_i')) =
        \begin{pmatrix}f_1\\\vdots\\f_k\end{pmatrix}
        (a_1'\cdot ...\cdot a_k') = \begin{pmatrix}f_1\\\vdots\\f_k\end{pmatrix}
        (a_1\cdot ...\cdot a_k)P^{-1} = PP^{-1} = E
        \Longrightarrow\\\\ a_1',...,a_k'$ -- ЛНЗ
        $\Longrightarrow a_1,..., a_k$ -- ЛНЗ.

    \end{proof}
    \newpage
    \begin{center}
        \begin{Large}
            \textbf{Глава II. Линейные отображения и операторы}
            
            \textbf{\S1. Линейные отображения, их матрицы.
            Ядро и образ.}
        \end{Large}
    \end{center}
    Пусть $V, V'$ -- векторные пространства над полем $F$.
    \begin{definition}
        $\varphi:\ V \longmapsto V'$ -- \textit{линейное отображение},
        если:\\1. $\forall v_1, v_2 \in V:\ \varphi(v_1 + v_2) =
        \varphi(v_1) + \varphi(v_2)$\\2. $\forall v \in V,
        \lambda \in F: \varphi(\lambda v) = \lambda \varphi(v)$
    \end{definition}
    \begin{consequense}
        $\varphi(0_V) = 0_{V'}$
    \end{consequense}$
    \varphi: V \longmapsto V$ -- \textit{линейный оператор} на $V$.
    \begin{definition}
        \textit{Ядро линейного отображения} $\varphi: V
        \longmapsto V'\ (V \longmapsto V).$
    \end{definition}
    $Ker \varphi = \{v \in V|\ \varphi(v) = 0_{V'}\}$
    (в частности, $0_V \in Ker\varphi$).

    Образ линейного отображения $\varphi$ (или множество
    значений $\varphi$) -- это\\ $Im\varphi (=\varphi(v)) = 
    \{v' \in V'|\ \exists v \in V: \varphi(v) = v'\}.$

    \begin{subtheorem}
        1. $Ker\varphi$ - подпространство в $V$.

        \ \ \ \ \ \ \ \ \ \ \ \ \ \ \ \ \ \ \ \ 
        2. $Im\varphi$ - подпространство в $V'$.
    \end{subtheorem}
    \begin{proof}
        $\\1.\ 0_V \in Ker\varphi$, если $\varphi(v_1) =
        \varphi(v_2) = 0_{V'}$, то $\varphi(v_1 + v_2) = 
        \varphi(v_1) + \varphi(v_2) = 0_{V'}$, если $\lambda
        \in F$, то $\varphi(\lambda v_1) = \lambda \varphi(v_1)
        = 0_{V'}.$\\
        $2.\ 0_{V'} \in Im\varphi$, если $v_1' = \varphi(v_1),\
        v_2' = \varphi(v_2) \Longrightarrow \forall \lambda
        \in F,\ \lambda v_1' = \varphi(\lambda v_1) \in Im\varphi$

    \end{proof}
    \begin{theorem}
        Если $dimV < \infty$, то $dimIm\varphi = dimV - dimKer
        \varphi$
    \end{theorem}
    \begin{proof}
        Если фикс. базис $e_1,..., e_n$ в $V$, то
        $Im\varphi = \langle \varphi(e_1),..., \varphi(e_n) \rangle$
        $\forall v \in V:\ v = \sum\limits_{i=1}^nx_ie_i
        \Longrightarrow \varphi(v) = \sum\limits_{i=1}^n
        x_i\varphi(e_i).$ В частности, $dimIm\varphi = m \leq n.$
        Выберем базис в $Im\varphi:\ f_1,...,f_m.\ \forall
        j = 1,...,m\ \exists a_j \in V:\ \varphi(a_j) = f_j$.
        \\Поймем, что векторы $a_1,...,a_m$ -- ЛНЗ.
        $\exists \alpha_1,...,\alpha_m \in F$\\
        Допустим, что $\alpha_1a_1,...,\alpha_ma_m = 0
        \Longrightarrow \varphi(\sum\limits_{j=1}^m\alpha_j
        a_j) = \sum\limits_{j=1}^m\alpha_j\varphi(a_j) =
        \sum\limits_{j=1}^m\alpha_jf_j = 0_{V'}$
        Т.к. $f_j$ -- ЛНЗ $\Longrightarrow\alpha_j = 0, \ 
        \forall_j = 1,...,m$\\Возьмем произвольный $v' \in Im
        \varphi:\ v' = \sum\limits_{j=1}^m\beta_jf_j = 0_{V'}$\\
        По определению образа $\exists v \in V:\ \varphi(v) = v'$\\
        Тогда рассмотрим $\varphi(v - \sum\limits_{j=1}^m\beta_j
        a_j) = v' - \sum\limits_{j=1}^m\beta_jf_j = 0_{V'}
        \Longrightarrow v - \sum\limits_{j=1}^m\beta_ja_j \in
        Ker\varphi$\\Выберем базис в $Ker\varphi:\
        \{b_1,..., b_r\}\Longrightarrow$ $$v - \sum\limits_{j=1}^m
        \beta_ja_j = \sum\limits_{k=1}^r\gamma_kb_k \Longrightarrow
        v = \sum\limits_{j=1}^m\beta_ja_j + \sum\limits_{k=1}^r
        x_kb_k$$\\Если взять $v$ произвольным, $v' = \varphi
        (v)$ -- разлагается по базису $f_1,...,f_m$ в $Im\varphi$,
        то предыдущие выкладки остаются в силе $\Longrightarrow
        \forall v \in V$ -- линейная комбинация векторов
        $\{a_j;b_k\}$. Эти векторы ЛНЗ в $V \Longrightarrow$
        базис в $V \Longrightarrow dimV = m+r.$\\Допустим,
        что $$\exists \beta_j\gamma_j:\ \sum\limits_{j=1}^m
        \beta_ja_j + \sum\limits_{k=1}^r\gamma_kb_k = 0_V
        \Longrightarrow \varphi(\sum\limits_j + \sum\limits_k)
        = \sum\limits_{j=1}^m\beta_jf_j + 0 = 0_{V'}
        \Longrightarrow$$ $$b_j = 0,\ j=1,...,m\Longrightarrow
        \undermat{\text{ЛНЗ}}{\\\sum\limits_{k=1}^r\gamma_kb_k}
        = 0\Longrightarrow \gamma_k = 0,\ k=1,...,r.$$
    \end{proof}
    
    Пусть $e = {e_1,...,e_n}$ -- базис в $V,\ \forall \in V:\ 
    v = \sum\limits_{j=1}^nx_je_j\Longrightarrow \varphi(v) =
    \sum\limits_{j=1}^nx_j\varphi(e_j)$.\\
    Если в $V'$ задан базис $f = \{f_1,...,f_m\}$ и известно
    разложение векторов $\varphi(e_j)$ по этому базису,\ 
    $\varphi(e_j) = \sum\limits_{i=1}^ma_{ij}f_i$, то можно
    вычислить $\varphi(v)$.\\$A_{\varphi,e,f} = (a_{ij})$ --
    матрица линейного отображения $\varphi$ в паре базисов
    $e$ и $f$. Таким образом, столбцы матрицы $A_\varphi$ --
    столбцы координат $\varphi(e_1),...,\varphi(e_m)$ в базисе
    $f$. Для $\varphi:\ V\longmapsto V$ по умолчанию $f=e$
    (второй базис равен первому), остается $A_{\varphi,e}$ --
    матрица линейного оператора $\varphi$ в базисе $e\\$.

    $V \mapsto V^{**} = (V^*)^*$ -- канонический изоморфизм.\\
    $\langle f|v \rangle$ -- значение функционала $f$ на векторе $v$\\
    $\langle f|$ -- bra-vector\\ $|v \rangle$ -- ket-vector.

    \begin{theorem}
        $\\$1. Если $dimV < \infty,\ \varphi: V \mapsto V'$ -- 
        линейное отображение, то $dimIm\varphi = dimV - dimKer
        \varphi$\\
        2. $Im\varphi \cong V/Ker\varphi$
    \end{theorem}
    \begin{proof}
        $\\1.$ Пусть $dimIm\varphi = m \leq dimV'$, выберем в образе
        $Im\varphi$\ \ базис $f_1,...,f_m.\\\forall j,\ 1\leq j\leq m,\ \exists\ 
        e_j \in V: \varphi(e_j) = f_j.$ $$\forall v \in V: v' = 
        \varphi(v) = \sum\limits_{j=1}^m\lambda_jf_j = 
        \varphi(\sum\limits_{j=1}^m\lambda_je_j) \Longrightarrow
        \varphi(v - v_1), \text{ т.е. } v - v_1 \in Ker\varphi.$$
        Выберем в $Ker\varphi$ базис: $c_1,..., c_r\ (z = dimKer\varphi)
        \Longrightarrow v - v_1 = \sum\limits_{k=1}^r\mu_kc_k
        \Longrightarrow\\ v = \sum\limits_{j=1}^m\lambda_je_j + 
        \sum\limits_{k=1}^r\mu_kc_k \Longrightarrow V = \langle 
        e_1,...,e_m, c_1,...,c_r \rangle$\\
        $e_1,...,e_m, c_1,...,c_r$ -- ЛНЗ, если: $\exists\ \varepsilon_j,
        \gamma_k \in F: \varphi(\sum\limits_{j=1}^m\varepsilon_je_j + 
        \sum\limits_{k=1}^k\gamma_ke_k) = 0_V \Longrightarrow
        \sum\limits_{j=1}^m\undermat{\text{ЛНЗ}}{\varepsilon_jf_j} = 0_{V'} \Longrightarrow
        \varepsilon_j = 0,\ 1\leq j\leq m \Longrightarrow 
        \sum\limits_{k=1}^r\undermat{\text{ЛНЗ}}{\gamma_ke_k} = 0 \Longrightarrow 
        \gamma_k = 0,\ 1\leq k\leq r.$
        $\\\\$2. \%Rem: Если $U \subseteq V$, то $V/U = \overline{v} \equiv
        \{v + u|\ u \in U\} = v + U$ -- факторпространство $V$
        по $U$, где $Ker\varphi := U$\\Рассмотрим отображение $\pi: V/U \mapsto
        Im\varphi \subseteq V',\ \pi(\overline{v}) = \varphi(v).$\\
        Корректность определения: если $v_1 \in V:
        \overline{v_1} = \overline{v}$, то $v_1 = v + u_1,
        \\\text{ где } 
        u_1 \in U = Ker\varphi \Longrightarrow \varphi(v_1) =
        \varphi(v) + \undermat{\equiv 0}{\varphi(u_1)} \Longrightarrow
        \pi(\overline{v_1}) = \pi(\overline{v})$

        $\pi$ -- линейное отображение:
        $$\pi(\overline{v_1} + \overline{v_2}) = 
        \pi(\overline{v_1 + v_2}) = \varphi(v_1 + v_2) = 
        \varphi(v_1) + \varphi(v_2) = \pi(\overline{v_1}) + 
        \pi(\overline{v_2}),\ \forall v \in V, \lambda \in F.$$
        $\pi$ -- биективное:\\
        Сюръективность: $\forall v' \in Im\varphi\ \exists\ v
        \in V: \varphi(v) = v' \Longrightarrow \pi(\overline{v}) =
        \varphi(v) = v'$.\\
        Инъективность: допустим, что $\ \pi(\overline{v_1}) =
        \pi(\overline{v_2}) \Longrightarrow \varphi(v_1) =
        \varphi(v_2) \Longrightarrow\\ v_2 - v_1 \in
        Ker\varphi \Longrightarrow \exists\ u \in U = Ker
        \varphi: v_2 = v_1 + u \Longrightarrow \overline{v_2} =
        \overline{v_1}$

    \end{proof}
    \underline{\textbf{Матрицы линейного отображения}} $\varphi: V \mapsto V'$\\
    $\forall v = \sum\limits_{j=1}^nx_je_j.\ (e = 
    (e_1,...,e_n))$ -- базис в $V$,\ 
    $f = (f_1,...,f_m)$ -- базис в $V',\\(dimV = n,\ dimV' = m)$
    $\Longrightarrow$ $$\varphi(v) = \sum\limits_{j=1}^nx_j
    \varphi(e_j) = \sum\limits_{j=1}^nx_j\sum\limits_{i=1}^ma_{ij}f_i = 
    \sum\limits_{i=1}^m(\sum\limits_{j=1}^na_{ij}x_j)f_i\ \ \ \ (1)$$
    Если $\varphi(e_j) = \sum\limits_{i=1}^ma_{ij}f_i,\ 
    (a_{ij}) = A_{\varphi,e,f}.$\\
    \underline{Обозначение:} $y = \varphi(v),\ v = eX_e^\uparrow \Longrightarrow
    y = \varphi(v) = fY_f^\uparrow = \sum\limits_{i=1}^my_if_i.$\\
    $(1) \Longrightarrow Y_e^\uparrow = A_{\varphi,e ,f}X_e^\uparrow\ \ \ (2)$\\
    $Ker\varphi = \{v = eX_e^\uparrow|\varphi(v) = 0\}$.
    Согласно формуле (2), $v \in Ker\varphi \Longleftrightarrow\\
    A_{\varphi,e,f}X_e^\uparrow = 0$ -- ОСЛУ с матрицей $A_\varphi
    \Longrightarrow dimKer\varphi = n - rkA_\varphi$\\
    Но $Im\varphi = \langle \varphi(e_1),...,\varphi(e_n) \rangle$ -- линейная
    оболочка столбцов матрицы $A_\varphi \Longrightarrow
    dimIm\varphi = rkA_\varphi$. Получили матричное доказательство
    формулы для $dimIm\varphi$.

    \begin{lemma}
        Линейное отображение $\varphi: V \mapsto V'$ инъективно
        $\Longleftrightarrow Ker\varphi = \{0\}\\
        (\text{и }\Longrightarrow \varphi(V) \cong V)$
    \end{lemma}
    \begin{proof}
        $\\\Longrightarrow:\ \varphi$ - инъективно. Знаем, что
        $\varphi(0_V) = 0_{V'}$, а т.к. $\varphi$ инъективно,
        то $0_V$, единственный вектор из $V,\ \mapsto
        0_{V'} \Longrightarrow Ker\varphi = \{0_V\}$.
        $\\\Longleftarrow:$ Пусть $\varphi(v_1) = \varphi(v_2)
        \Longrightarrow \varphi(v_1 - v_2) = 0_{V'}\Longrightarrow
        v_1 - v_2 \in Ker\varphi = \{0\} \Longrightarrow
        v_1 = v_2$

    \end{proof}

    \textbf{Задача:} Пусть $\varphi:\ V \mapsto V',\ A_{\varphi,e,f}$ --
    матрица $\varphi$. При каких усл. на $A_\varphi:\\
    1.\varphi$ инъективно?\\
    $2.\varphi$ сюръективно?\\
    $3.\varphi$ биективно?

    \textbf{Примеры:}\\1. $V = \mathbb{R}_n[x] = \{a_0 + a_1x +...+ a_nx^n\ |
    \ a_i \in \mathbb{R}\}$ -- пространство многочленов степени $\leq n$.
    Рассмотрим базис: $1, \frac{x}{1!},...,\frac{x^n}{n!}$
    $$\frac{x^2}{2!} \mapsto \frac{x}{1!};\ \varphi(e_j) = e_{j-1},
    \ 1\leq j \leq n.$$ $$\varphi = \frac{d}{dx},\ \frac{d}{dx}:
    V \mapsto V\ \ (\text{либо} V' = \mathbb{R}_{n-1}[x])$$
    $Ker\varphi = \{const\},\ Im\varphi = \mathbb{R}_{n-1}[x]$
    $$A_{\varphi,e} = 
    \begin{pmatrix}
        0 & 1 & 0 & \dots & 0\\
        0 & 0 & 1 & \dots & 0\\
        \vdots & \vdots & \vdots & \ddots & \vdots\\
        0 & 0 & 0 & \dots & 1\\
        0 & 0 & 0 & \dots & 0
    \end{pmatrix}$$
    Жорданова клетка порядка $n+1$ с диагональным элементом $\lambda = 0$.\\
    2. Пусть $V = U_1 \oplus U_2$ -- прямая сумма подпространств.\\
    Определим $\varphi: V \mapsto V,\ \forall v = u_1 + u_2,\ 
    \varphi(v) = u_1$ -- проектирование $V$ на подпространство
    $U_1\ ||\ U_2$.\\
    $\varphi$ -- линейный оператор.
    $Ker\varphi = U_2,\ Im\varphi = U_1$\\
    Если $dim\ U_1 = k,\ dim\ U_2 = n - k$.\\
    Рассмотрим $\undermat{\text{базис } U_1}{e_1,...,e_k},
    \undermat{\text{базис } U_2}{e_{k+1},...,e_n} \Longrightarrow
    A_{\varphi,e} = 
    \begin{pmatrix}
        E_k & \vline & 0\\
        \hline 0 & \vline & 0
    \end{pmatrix}$
    \newpage
    \underline{\textbf{Вычисление матрицы}} $A_{\varphi,e,f}:$\\
    $A_\varphi e_j^\uparrow$ -- j-й столбец матрицы $A_\varphi$
    в базисе $f$, где $1 \leq j \leq n.$\\
    \textbf{Рассмотрим задачу:} даны ЛНЗ векторы $a_1,..., a_n$ в $V$ и
    некоторые векторы $b_1,...,b_n$ в $V'$. Найти матрицу
    линейного отображения $\varphi: V \mapsto V'$ такого, чтобы
    $\varphi(a_i) = b_i,\ 1\leq i \leq n$.\\ По формуле,
    $b_j^\uparrow = A_\varphi a_j^\uparrow,\ 1 \leq j \leq n
    \Longleftrightarrow A_\varphi(a_1^\uparrow,...,a_n^\uparrow) =
    (b_1^\uparrow,...,b_n^\uparrow)$\\ Уравнение вида: $XA = B,\ 
    detA \neq 0 \Longrightarrow X = BA^{-1}$\\
    $\\\begin{pmatrix}A\\\hline B\end{pmatrix}
    \underset{\text{столбцов}}{\overset{\text{Э.П.}}
    {\sim}}\begin{pmatrix}E\\\hline X\end{pmatrix}
    \Longrightarrow X = A_{\varphi,e,f}$\\
    $\\(A^T|B^T) \underset{\text{строк}}{\overset{\text{Э.П.}}{\sim}} (E|X^T)$
    
    \underline{\textbf{Изменение матрицы линейного
    отображения (оператора) при}}\\
    \underline{{\textbf{замене базисов:}}}

    Пусть $e, e':\ e' = eC_{e \to e'}$ в $V$ и
    $f, f':\ f' = fD_{f \to f'}$ в $V'$.

    Но $X_{e'}:\ X_e = C_{e \to e'}X_{e'} \ \ \ (3)\tab[5cm]
    D_{f \to f'}Y_{f'} = A_{\varphi,e,f}C_{e \to e'}X_{e'}$

    \ \ \ \ \ 
    $Y_{f'}:\ Y_f = D_{f \to f'}Y_{f'}\ \ \ (3')\tab[5cm]
    Y_{f'} = (D^{-1}A_\varphi C)X_{e'} \ \ \ (*)$

    $(*)$ Сравним это с формулой:
    $$Y_{f'} = A_{\varphi,e',f'},\ \forall x\in F^n, Y\in F^m
    \Longrightarrow A_{\varphi,e',f'} = D_{f\to f'}^{-1}
    A_{\varphi,e,f} C_{e \to e'}$$

    В качестве $X$ по очереди взять столбцы $E_n$, в качестве
    $Y$ -- столбцы $E_m.$
    
    $\\$Если $\varphi: V \longmapsto V$ -- линейный оператор,
    $e, e'$ -- базисы, $C = C_{e \to e'}$, то\\ 
    $a_{\varphi, e'} = C^{-1}A{\varphi,e}\ C$\ $(*)$

    $\\$\underline{Термин:} Две матрицы называются \textit{подобными},
    если $$\exists C:\ A' = C_{-1}AC\ (detC \neq 0)$$

    \begin{subtheorem}
        Если $A$ и $A'$ подобные, то $|A'| = |A|$ и $rkA'=rkA.$  
    \end{subtheorem}
    
    (*) $\Longleftrightarrow CA' = AC$. Составим матрицу $(C|AC)
    \underset{\text{строк}}{\overset{\text{Э.П.}}{\sim}}  (E|A')$.
    \newpage
    \begin{center}
        \begin{Large}    
            \textbf{\S2. Действия над линейными отображениями
             и операторами.}
        \end{Large}
    \end{center}    

    $\varphi, \psi: V \longmapsto V'$ --линейное.\\
    $\forall v \in V\ (\varphi + \psi)(v) := \varphi(v) + \psi(v)$\\
    $\forall \lambda \in F\ (\lambda \varphi)(v) := \lambda \varphi(v)$   
    \begin{subtheorem}
        $\\1. \mathcal{L}(V,V')$ является векторным пространством над $F$\\
        2. Если $dimV = n,\ dimV' = m$, то $\mathcal{L} (V,V') \cong
        M_{m\times n}(F)$  
    \end{subtheorem}    
    \begin{proof}
        1. "Очевидно".

        2. Если фиксировать базисы $e$ в $V$,\ $f$ в $V'$,
        то $\forall \varphi \longleftrightarrow  A_{\varphi,e,f},$
        причем:\\ $A_{\varphi + \psi} = A_\varphi + A_\psi$,\ \
        $A_{\lambda \varphi} = \lambda A_\varphi$
        
        Если $X$ -- столбец координат вектора $v$,\ $v' = \varphi(v)$,
        $Y$ -- столбец координат $v'$.
        
        $Y = A_\varphi X$,
        
        $\widetilde{Y} $ -- столбец координат вектора $\psi(v)$, то
        $\widetilde{Y} = A_\psi X,\ (\varphi + \psi)(v) =\\ \varphi(v) + \psi(v)
        \Longrightarrow \varphi(v) + \psi(v)$ имеет столбец
        координат $Y + \widetilde{Y} = A_\varphi X + A_\psi X =\\
        (A_\varphi + A_\psi)X = A_{\varphi + \psi}X \Longrightarrow 
        \varphi \longmapsto A_{\varphi,e,f}$ -- изоморфизм ($dim
        \mathcal{L}(V,V') = n\cdot m$)

    \end{proof}

    \underline{\textbf{Умножение (композиция) линейных отображений (операторов):}}

    Пусть $V \overset{\psi}{\longmapsto} V' \overset{\varphi}{\longmapsto}  V''$ 

    \begin{subtheorem}
        Если $\varphi,\ \psi$ -- линейные, то $\varphi\cdot \psi$ --
        линейное  
    \end{subtheorem}

    В частности, для линейных операторов: $\varphi\cdot \psi$ --
    линейный оператор в этом же пространстве.
    
    \underline{Обозначим:} $\mathcal{L}(V)$ -- множество всех линейных операторов на $V$.
    
    \begin{theorem}
        С операциями $+,\ \lambda\ \cdot,$ и $\cdot$\ $\mathcal{L}(V)$ является
        [линейной] алгеброй; если $dimV = n$, то $L(V) \cong M_n(F)$  
    \end{theorem}
    
    ($A$ -- линейная алгебра, если на $A$ заданы операции
    $+, \cdot, \lambda\ \cdot$, такие что
    \begin{enumerate}
        \item $A_+ $\text{ -- кольцо ассоциативное}
        \item $A_{+,\lambda\cdot,\ \cdot}$
        \text{ -- векторное пространство}
        \item $\lambda (ab) = (\lambda a)b, \forall \lambda
        \in F,\ a,b \in A.$)
    \end{enumerate}
    
    \begin{proof}
        1, 2 уже проверены.\\
        3. $\forall v \in V,\ (\lambda(ab))(v) = \lambda
        ((ab)(v)) = \lambda a(b(v)) = (\lambda a)(b(v))
        \Longrightarrow \lambda(ab) = (\lambda a)b$, для 
        $a,b \in \mathcal{L}(V)$.  
    \end{proof}

    \textbf{Геометрическое доказательство неравенства} $rk(AB) \leq
    \begin{cases}
        rkA\\rkB
    \end{cases}$\\ (для любых матриц $A,B$, т.ч. $\exists AB$).
    \begin{proof}
        Если $A_{m\times n},\ B_{n\times p}$, то $\mathbb{R}^p
        \overset{B}{\longmapsto} \mathbb{R}^n \overset{A}{\longmapsto}
        \mathbb{R}^m$\\
        $rk(AB) = dim(AB(\mathbb{R}^p)) \leq dimB(\mathbb{R}^p)\\
        A(\mathbb{R}^n) \supseteq (AB)(\mathbb{R}^p)
        \Longrightarrow rk(AB) = dimIm(AB) \leq dimImA = rkA.$

        Случай равенства. Если $|A| \neq 0,$ то $A$ задает
        изоморфизм между $\mathbb{R}^n$ и $\mathbb{R}^m
        \Longrightarrow dimB(\mathbb{R}^p) = dim(AB(\mathbb{R}^p))$,
    т.е. $rkB = rkAB$

    \end{proof}
    \textbf{Многочлены от линейного оператора:}

    Пусть $p(t) = a_0t^m + a_1t_{m-1}+...+a_{m-1}t + a_m,\ 
    \forall a_i \in F,\ i = 0,...,m$.\\ Тогда $p(\varphi) = 
    a_0\varphi^m + a_1\varphi^{m-1}+...+a_{m-1}\varphi + a_m \varepsilon$
    
    $\varepsilon(v) = v, \forall v \in V$ -- тождественный оператор.
    
    Многочлен $p(t)$ -- \textit{аннулирующий многочлен оператора} $\varphi$,
    если $p(\varphi) = 0$ (нулевой оператор).
    \begin{theorem}
        Для любого оператора $\varphi: V_n \longmapsto V_n$
        существует аннулирующий многочлен степени $\leq n^2-1$. 
    \end{theorem}
    \begin{proof}
        Т.к. $dim \mathcal{L}(V) = n^2$, то операторы, $\varphi^0 = \varepsilon,
        \varphi, \varphi^2,..., \varphi^{n^2}$ -- ЛЗ. $\Longrightarrow
        \exists a_0, a_1,..., a_{n^2} \in F,\ 
        a_0\varepsilon + a_1\varphi +...+a_{n^2}\varphi^{n^2} = 0
        \Longrightarrow p(t) = a_0 + a_1t +...+ a_{n^2}t^n$ --
        аннулирующий для $\varphi$.
        
    \end{proof}

    \begin{center}
        \begin{Large}    
            \textbf{\S3. Инвариантные пространства. Собственные
            векторы.}
        \end{Large}
    \end{center}

    Пусть $\varphi: V \longmapsto V$ -- линейный оператор.
    
    \begin{definition}
        Подпространство $U \subseteq V$ называется \textit{
        инвариантым} относительно $\varphi$ (или для $\varphi$,
        или $\varphi$-инвариантным), если $\forall u \in U
        \Longrightarrow \varphi(u) \in U$, т.е. $\varphi(U)
        \subseteq U$. 
    \end{definition}

    \begin{definition}
        \textit{Ограничение (сужение) оператора} $\varphi$
        на инвариантное подмножество $U$ -- это оператор. 
    \end{definition}

    \underline{Обозначение:} $\varphi|_v: U \longmapsto U$, а именно,
    $\forall u \in U: \varphi|_u(U) = \varphi(U)$
    \newpage
    \textbf{Примеры:}
    \begin{enumerate}
        \item $\varphi = \frac{d}{dx}: \mathbb{R}[x]
        \longmapsto \mathbb{R}[x]$,\ \ $p(x) \longmapsto p'(x)$
        
        Для $\forall n = 0,1,...,\ \mathbb{R}_n[x] = \{p(x)|
        deg\ p(x)\leq n\}$ являются инвариантными подпространствами

        Вопрос: есть ли другие инвариантные подпространства?
        \item $\varphi$ -- проектирование (проектор) $V
        \longmapsto V = U_1 \oplus U_2$
        
        $\varphi(u_1 + u_2) = u$ на $U_1$ вдоль $U_2$.
        
        Инвариантные: $U_1, U_2$, а также $U_1' \oplus U_2'$,
        для $\forall\ U_1' \subseteq U_1,\ U_2' \subseteq U_2$.
    \end{enumerate}
    
    \begin{theorem}
        Если $\varphi: V \longmapsto  V$, то $\varphi$-инвариантными
        являются:
        \begin{itemize}
            \item $Ker \varphi$
            \item $Im \varphi$
            \item Любое пространство $U$, $Im \varphi \subseteq 
            U \subseteq V$
            \item $\forall k = 1,2,...,\ Ker \varphi^k$
            \item $Im \varphi^k$ 
        \end{itemize}
    \end{theorem}

    \begin{theorem}
        Если $U_1,..., U_k$ -- $\varphi$-инвариантные, то $\varphi$-инвариантные:
        \begin{itemize}
            \item $U_1+...+U_k$
            \item $U_1 \cap...\cap U_k$ 
        \end{itemize}        
    \end{theorem}

    \begin{lemma}
        (Вид матрицы $A_\varphi$ при наличии инвариантов 
        подпространства)
        \begin{enumerate}
            \item Пусть $U_1$ -- $\varphi$-инвариантное
            пространство, $\{0\} \neq U_1 \neq V$.\\ Тогда
            в $V\ \exists$ базис, в котором\\ $A_\varphi = 
            \begin{pmatrix}
                \underset{n\times m}{B} & \underset{m\times n-m}{D}\\
                \underset{(n-m)\times n}{0} & \underset{(n-m)\times (n-m)}{C} 
            \end{pmatrix} \text{, причем } B = A_{\varphi|U_1},\ \  dimU_1 = m$ 
            \item Пусть $V = U_1 \oplus U_2$, причем $U_1, U_2$ --
            ненулевые инвариантные подпространства, тогда в
            $V\ \exists$ базис, в котором\\ $A_\varphi = 
            \begin{pmatrix}
                B & 0\\
                0 & C 
            \end{pmatrix}$, причем $C = A_{\varphi|U_2}$ 
        \end{enumerate}
    \end{lemma}
    \begin{proof}
        $\\$1. Надо взять $e_1,..., e_m$ -- базис в $U_1$,
        дополнить его произвольно до базиса в $V$.\\
        Тогда для $j=1,...,m,\ \varphi(e_j) \in U_1,\ 
        \varphi(e_j) = \sum\limits_{i=1}^{m} b_ie_i$\\
        2. Если базис $e_1,..., e_m$ пространства
        $U_1$ объединить с базисом $e_{m+1},...,
        e_n$\\ пространства $U_2$, то в полученном базисе
        $A_\varphi = 
    \begin{pmatrix}
        B & 0\\
        0 & C 
    \end{pmatrix}$
    
    \end{proof}
    \begin{subtheorem}
        Верное и обратное
    \end{subtheorem}
    




























\end{document}